% GENERATED FILE. Edit canonical lessons, then run npm run generate:books.
% canonical-source-hash: fnv1a64:6df6bf40791ad3cf
% canonical-lessons: UR-C382-inam, UR-C382-ezaz, UR-C382-tamgha, UR-C382-mustahiq, UR-C382-sanvarna

\chapter{Prizes and Rewards}
\label{ch:ur-prizes-and-rewards}
\hlchaptermodality{382}

\begin{chapteropening}
\textbf{By the end of this chapter:} \emph{I can say what someone will win if they do well.}

Complete the last lesson of chapter 382: I can say what someone will win if they do well.
\end{chapteropening}

\section[\texorpdfstring{\ur{انعام} (inām)}{انعام (inām)}]{\ur{انعام} (inām) — a prize; a reward}

\label{lesson:UR-C382-inam}



\begin{quote}
Before the new one: say the Urdu for a marvel, then the Urdu for well done.
\end{quote}

\subsection*{You'll want to know: \ur{انعام}}
\textbf{\ur{انعام}} — \emph{inām} — \textquotedblleft{}a prize; a reward\textquotedblright{}.

\begin{quote}
\ur{اگر} \ur{تم} \ur{جیتے} \ur{تو} \ur{تمہیں} \ur{انعام} \ur{ملے} \ur{گا۔}
\end{quote}

\begin{quote}
\emph{agar tum jīte to tumheṅ inām mile gā.}
\end{quote}

\begin{quote}
If you win, you will get a prize.
\end{quote}

Masculine.

\begin{cousinweb}
\textbf{\ur{انعام}} (\emph{inām}) is usually traced, through Persian, to Arabic \emph{inʿām}, \textquotedblleft{}a bestowing of favour\textquotedblright{}.
\end{cousinweb}

\begin{grammarlens}[title={What you've built}]
One more word for saying what happens if.
\end{grammarlens}

\subsection*{Your turn}
\begin{itemize}
\raggedright
  \item \emph{Say it:} \emph{inām}
  \item \emph{Say it:} \emph{inām}, once more
  \item \emph{Say it:} \emph{agar tum jīte to tumheṅ inām mile gā.} — If you win, you will get a prize.
  \item \emph{Recall:} say the Urdu for a marvel, then the Urdu for well done, then say \emph{inām} again
\end{itemize}

\begin{tcolorbox}[breakable,colback=teal!4,colframe=teal!35!black,title={Before you move on}]
What does \ur{انعام} mean? (\textquotedblleft{}A prize; a reward\textquotedblright{}.) Say it once more.
\end{tcolorbox}
\section[\texorpdfstring{\ur{اعزاز} (ezāz)}{اعزاز (ezāz)}]{\ur{اعزاز} (ezāz) — an honour; an award}

\label{lesson:UR-C382-ezaz}



\begin{quote}
Before the new one: say the Urdu for well done, then the Urdu for a prize.
\end{quote}

\subsection*{You'll want to know: \ur{اعزاز}}
\textbf{\ur{اعزاز}} — \emph{ezāz} — \textquotedblleft{}an honour; an award\textquotedblright{}.

\begin{quote}
\ur{اگر} \ur{اسے} \ur{یہ} \ur{اعزاز} \ur{ملا} \ur{تو} \ur{سب} \ur{خوش} \ur{ہوں} \ur{گے۔}
\end{quote}

\begin{quote}
\emph{agar use yih ezāz milā to sab khush hoṅ ge.}
\end{quote}

\begin{quote}
If she receives this award, everyone will be happy.
\end{quote}

Masculine.

\begin{cousinweb}
\textbf{\ur{اعزاز}} (\emph{ezāz}) is usually traced, through Persian, to Arabic \emph{iʿzāz}, \textquotedblleft{}honouring\textquotedblright{}, of the same family as \textbf{\ur{عزت}}, \textquotedblleft{}honour\textquotedblright{}.
\end{cousinweb}

\begin{grammarlens}[title={What you've built}]
One more word for saying what happens if.
\end{grammarlens}

\subsection*{Your turn}
\begin{itemize}
\raggedright
  \item \emph{Say it:} \emph{ezāz}
  \item \emph{Say it:} \emph{ezāz}, once more
  \item \emph{Say it:} \emph{agar use yih ezāz milā to sab khush hoṅ ge.} — If she receives this award, everyone will be happy.
  \item \emph{Recall:} say the Urdu for well done, then the Urdu for a prize, then say \emph{ezāz} again
\end{itemize}

\begin{tcolorbox}[breakable,colback=teal!4,colframe=teal!35!black,title={Before you move on}]
What does \ur{اعزاز} mean? (\textquotedblleft{}An honour; an award\textquotedblright{}.) Say it once more.
\end{tcolorbox}
\section[\texorpdfstring{\ur{تمغہ} (tamghā)}{تمغہ (tamghā)}]{\ur{تمغہ} (tamghā) — a medal}

\label{lesson:UR-C382-tamgha}



\begin{quote}
Before the new one: say the Urdu for a prize, then the Urdu for an award.
\end{quote}

\subsection*{You'll want to know: \ur{تمغہ}}
\textbf{\ur{تمغہ}} — \emph{tamghā} — \textquotedblleft{}a medal\textquotedblright{}.

\begin{quote}
\ur{اگر} \ur{اس} \ur{نے} \ur{دوڑ} \ur{جیتی} \ur{تو} \ur{اسے} \ur{تمغہ} \ur{ملے} \ur{گا۔}
\end{quote}

\begin{quote}
\emph{agar us ne dauṛ jītī to use tamghā mile gā.}
\end{quote}

\begin{quote}
If she wins the race, she will get a medal.
\end{quote}

Masculine.

\begin{cousinweb}
\textbf{\ur{تمغہ}} (\emph{tamghā}) is usually traced, through Persian, to Turkish \emph{tamga}, \textquotedblleft{}a seal, a stamp\textquotedblright{}: a stamped token of honour.
\end{cousinweb}

\begin{grammarlens}[title={What you've built}]
One more word for saying what happens if.
\end{grammarlens}

\subsection*{Your turn}
\begin{itemize}
\raggedright
  \item \emph{Say it:} \emph{tamghā}
  \item \emph{Say it:} \emph{tamghā}, once more
  \item \emph{Say it:} \emph{agar us ne dauṛ jītī to use tamghā mile gā.} — If she wins the race, she will get a medal.
  \item \emph{Recall:} say the Urdu for a prize, then the Urdu for an award, then say \emph{tamghā} again
\end{itemize}

\begin{tcolorbox}[breakable,colback=teal!4,colframe=teal!35!black,title={Before you move on}]
What does \ur{تمغہ} mean? (\textquotedblleft{}A medal\textquotedblright{}.) Say it once more.
\end{tcolorbox}
\section[\texorpdfstring{\ur{مستحق} (mustahiq)}{مستحق (mustahiq)}]{\ur{مستحق} (mustahiq) — deserving}

\label{lesson:UR-C382-mustahiq}



\begin{quote}
Before the new one: say the Urdu for an award, then the Urdu for a medal.
\end{quote}

\subsection*{You'll want to know: \ur{مستحق}}
\textbf{\ur{مستحق}} — \emph{mustahiq} — \textquotedblleft{}deserving\textquotedblright{}.

\begin{quote}
\ur{اگر} \ur{وہ} \ur{مستحق} \ur{ہے} \ur{تو} \ur{اسے} \ur{انعام} \ur{ملنا} \ur{چاہیے۔}
\end{quote}

\begin{quote}
\emph{agar voh mustahiq hai to use inām milnā chāhiye.}
\end{quote}

\begin{quote}
If he deserves it, he should get the prize.
\end{quote}

\begin{cousinweb}
\textbf{\ur{مستحق}} (\emph{mustahiq}) is usually traced, through Persian, to Arabic \emph{mustaḥiqq}, \textquotedblleft{}deserving, having a right\textquotedblright{}, of the same family as \textbf{\ur{حق}}, \textquotedblleft{}a right\textquotedblright{}.
\end{cousinweb}

\begin{grammarlens}[title={What you've built}]
One more word for saying what happens if.
\end{grammarlens}

\subsection*{Your turn}
\begin{itemize}
\raggedright
  \item \emph{Say it:} \emph{mustahiq}
  \item \emph{Say it:} \emph{mustahiq}, once more
  \item \emph{Say it:} \emph{agar voh mustahiq hai to use inām milnā chāhiye.} — If he deserves it, he should get the prize.
  \item \emph{Recall:} say the Urdu for an award, then the Urdu for a medal, then say \emph{mustahiq} again
\end{itemize}

\begin{tcolorbox}[breakable,colback=teal!4,colframe=teal!35!black,title={Before you move on}]
What does \ur{مستحق} mean? (\textquotedblleft{}Deserving\textquotedblright{}.) Say it once more.
\end{tcolorbox}
\section[\texorpdfstring{\ur{سنوارنا} (sanvārnā)}{سنوارنا (sanvārnā)}]{\ur{سنوارنا} (sanvārnā) — to tidy; to make (something) look its best}

\label{lesson:UR-C382-sanvarna}



\begin{quote}
Before the new one: say the Urdu for a medal, then the Urdu for deserving.
\end{quote}

\subsection*{You'll want to know: \ur{سنوارنا}}
\textbf{\ur{سنوارنا}} — \emph{sanvārnā} — \textquotedblleft{}to tidy; to make (something) look its best\textquotedblright{}.

\begin{quote}
\ur{اگر} \ur{تم} \ur{نے} \ur{کمرہ} \ur{سنوارا} \ur{تو} \ur{امی} \ur{خوش} \ur{ہوں} \ur{گی۔}
\end{quote}

\begin{quote}
\emph{agar tum ne kamrā sanvārā to ammī khush hoṅ gī.}
\end{quote}

\begin{quote}
If you tidy the room, your mother will be pleased.
\end{quote}

Past, with \ur{نے} (he or she …): \ur{اس} \ur{نے} \ur{بال} \ur{سنوارے}, \emph{us ne bāl sanvāre}.

\begin{cousinweb}
\textbf{\ur{سنوارنا}} (\emph{sanvārnā}) is inherited: CDIAL lists the Hindi cousin \emph{sãwārnā}, \textquotedblleft{}to arrange, to deck\textquotedblright{}, under Sanskrit \emph{saṃvārayati} (CDIAL 13021).
\end{cousinweb}

\begin{grammarlens}[title={What you've built}]
One more word for saying what happens if.
\end{grammarlens}

\subsection*{Your turn}
\begin{itemize}
\raggedright
  \item \emph{Say it:} \emph{sanvārnā}
  \item \emph{Say it:} \emph{sanvārnā}, once more
  \item \emph{Say it:} \emph{agar tum ne kamrā sanvārā to ammī khush hoṅ gī.} — If you tidy the room, your mother will be pleased.
  \item \emph{Recall:} say the Urdu for a medal, then the Urdu for deserving, then say \emph{sanvārnā} again
\end{itemize}

\begin{tcolorbox}[breakable,colback=teal!4,colframe=teal!35!black,title={Before you move on}]
What does \ur{سنوارنا} mean? (\textquotedblleft{}To tidy; to make (something) look its best\textquotedblright{}.) Say it once more.
\end{tcolorbox}
